% Options for packages loaded elsewhere
\PassOptionsToPackage{unicode}{hyperref}
\PassOptionsToPackage{hyphens}{url}
\documentclass[
]{article}
\usepackage{xcolor}
\usepackage[margin=1in]{geometry}
\usepackage{amsmath,amssymb}
\setcounter{secnumdepth}{-\maxdimen} % remove section numbering
\usepackage{iftex}
\ifPDFTeX
  \usepackage[T1]{fontenc}
  \usepackage[utf8]{inputenc}
  \usepackage{textcomp} % provide euro and other symbols
\else % if luatex or xetex
  \usepackage{unicode-math} % this also loads fontspec
  \defaultfontfeatures{Scale=MatchLowercase}
  \defaultfontfeatures[\rmfamily]{Ligatures=TeX,Scale=1}
\fi
\usepackage{lmodern}
\ifPDFTeX\else
  % xetex/luatex font selection
\fi
% Use upquote if available, for straight quotes in verbatim environments
\IfFileExists{upquote.sty}{\usepackage{upquote}}{}
\IfFileExists{microtype.sty}{% use microtype if available
  \usepackage[]{microtype}
  \UseMicrotypeSet[protrusion]{basicmath} % disable protrusion for tt fonts
}{}
\makeatletter
\@ifundefined{KOMAClassName}{% if non-KOMA class
  \IfFileExists{parskip.sty}{%
    \usepackage{parskip}
  }{% else
    \setlength{\parindent}{0pt}
    \setlength{\parskip}{6pt plus 2pt minus 1pt}}
}{% if KOMA class
  \KOMAoptions{parskip=half}}
\makeatother
\usepackage{color}
\usepackage{fancyvrb}
\newcommand{\VerbBar}{|}
\newcommand{\VERB}{\Verb[commandchars=\\\{\}]}
\DefineVerbatimEnvironment{Highlighting}{Verbatim}{commandchars=\\\{\}}
% Add ',fontsize=\small' for more characters per line
\usepackage{framed}
\definecolor{shadecolor}{RGB}{248,248,248}
\newenvironment{Shaded}{\begin{snugshade}}{\end{snugshade}}
\newcommand{\AlertTok}[1]{\textcolor[rgb]{0.94,0.16,0.16}{#1}}
\newcommand{\AnnotationTok}[1]{\textcolor[rgb]{0.56,0.35,0.01}{\textbf{\textit{#1}}}}
\newcommand{\AttributeTok}[1]{\textcolor[rgb]{0.13,0.29,0.53}{#1}}
\newcommand{\BaseNTok}[1]{\textcolor[rgb]{0.00,0.00,0.81}{#1}}
\newcommand{\BuiltInTok}[1]{#1}
\newcommand{\CharTok}[1]{\textcolor[rgb]{0.31,0.60,0.02}{#1}}
\newcommand{\CommentTok}[1]{\textcolor[rgb]{0.56,0.35,0.01}{\textit{#1}}}
\newcommand{\CommentVarTok}[1]{\textcolor[rgb]{0.56,0.35,0.01}{\textbf{\textit{#1}}}}
\newcommand{\ConstantTok}[1]{\textcolor[rgb]{0.56,0.35,0.01}{#1}}
\newcommand{\ControlFlowTok}[1]{\textcolor[rgb]{0.13,0.29,0.53}{\textbf{#1}}}
\newcommand{\DataTypeTok}[1]{\textcolor[rgb]{0.13,0.29,0.53}{#1}}
\newcommand{\DecValTok}[1]{\textcolor[rgb]{0.00,0.00,0.81}{#1}}
\newcommand{\DocumentationTok}[1]{\textcolor[rgb]{0.56,0.35,0.01}{\textbf{\textit{#1}}}}
\newcommand{\ErrorTok}[1]{\textcolor[rgb]{0.64,0.00,0.00}{\textbf{#1}}}
\newcommand{\ExtensionTok}[1]{#1}
\newcommand{\FloatTok}[1]{\textcolor[rgb]{0.00,0.00,0.81}{#1}}
\newcommand{\FunctionTok}[1]{\textcolor[rgb]{0.13,0.29,0.53}{\textbf{#1}}}
\newcommand{\ImportTok}[1]{#1}
\newcommand{\InformationTok}[1]{\textcolor[rgb]{0.56,0.35,0.01}{\textbf{\textit{#1}}}}
\newcommand{\KeywordTok}[1]{\textcolor[rgb]{0.13,0.29,0.53}{\textbf{#1}}}
\newcommand{\NormalTok}[1]{#1}
\newcommand{\OperatorTok}[1]{\textcolor[rgb]{0.81,0.36,0.00}{\textbf{#1}}}
\newcommand{\OtherTok}[1]{\textcolor[rgb]{0.56,0.35,0.01}{#1}}
\newcommand{\PreprocessorTok}[1]{\textcolor[rgb]{0.56,0.35,0.01}{\textit{#1}}}
\newcommand{\RegionMarkerTok}[1]{#1}
\newcommand{\SpecialCharTok}[1]{\textcolor[rgb]{0.81,0.36,0.00}{\textbf{#1}}}
\newcommand{\SpecialStringTok}[1]{\textcolor[rgb]{0.31,0.60,0.02}{#1}}
\newcommand{\StringTok}[1]{\textcolor[rgb]{0.31,0.60,0.02}{#1}}
\newcommand{\VariableTok}[1]{\textcolor[rgb]{0.00,0.00,0.00}{#1}}
\newcommand{\VerbatimStringTok}[1]{\textcolor[rgb]{0.31,0.60,0.02}{#1}}
\newcommand{\WarningTok}[1]{\textcolor[rgb]{0.56,0.35,0.01}{\textbf{\textit{#1}}}}
\usepackage{graphicx}
\makeatletter
\newsavebox\pandoc@box
\newcommand*\pandocbounded[1]{% scales image to fit in text height/width
  \sbox\pandoc@box{#1}%
  \Gscale@div\@tempa{\textheight}{\dimexpr\ht\pandoc@box+\dp\pandoc@box\relax}%
  \Gscale@div\@tempb{\linewidth}{\wd\pandoc@box}%
  \ifdim\@tempb\p@<\@tempa\p@\let\@tempa\@tempb\fi% select the smaller of both
  \ifdim\@tempa\p@<\p@\scalebox{\@tempa}{\usebox\pandoc@box}%
  \else\usebox{\pandoc@box}%
  \fi%
}
% Set default figure placement to htbp
\def\fps@figure{htbp}
\makeatother
\setlength{\emergencystretch}{3em} % prevent overfull lines
\providecommand{\tightlist}{%
  \setlength{\itemsep}{0pt}\setlength{\parskip}{0pt}}
\usepackage{bookmark}
\IfFileExists{xurl.sty}{\usepackage{xurl}}{} % add URL line breaks if available
\urlstyle{same}
\hypersetup{
  pdftitle={Reproducible Research\_Peer Assignment 1},
  pdfauthor={ITguy666},
  hidelinks,
  pdfcreator={LaTeX via pandoc}}

\title{Reproducible Research\_Peer Assignment 1}
\author{ITguy666}
\date{2026-03-15}

\begin{document}
\maketitle

For this assignment, we will first download the data from a personal
activity monitoring device that collects data at 5 minute intervals
throughout the day. The data consists of two months of data from an
anonymous individual collected during months of October and November
2012 and include the number of steps taken in 5 minute intervals each
day.

The variables included in this dataset are:

\begin{enumerate}
\def\labelenumi{\arabic{enumi})}
\item
  steps: Number of steps taking in a 5-minute interval (missing values
  are coded as NA)
\item
  date: The date on which the measurement was taken in YYYY-MM-DD format
\item
  interval: Identifier for the 5-minute interval in which measurement
  was taken
\end{enumerate}

The dataset is stored in a comma-separated-value (CSV) file and there
are a total of 17,568 observations in this dataset.

\begin{Shaded}
\begin{Highlighting}[]
\NormalTok{knitr}\SpecialCharTok{::}\NormalTok{opts\_chunk}\SpecialCharTok{$}\FunctionTok{set}\NormalTok{(}\AttributeTok{echo =} \ConstantTok{TRUE}\NormalTok{,}\AttributeTok{warning =} \ConstantTok{FALSE}\NormalTok{, }\AttributeTok{message =} \ConstantTok{FALSE}\NormalTok{)}
\DocumentationTok{\#\#Loading libraries and preprocessing the activity data}
\FunctionTok{library}\NormalTok{(ggplot2)}
\FunctionTok{library}\NormalTok{(lubridate)}
\end{Highlighting}
\end{Shaded}

\begin{verbatim}
## 
## Attaching package: 'lubridate'
\end{verbatim}

\begin{verbatim}
## The following objects are masked from 'package:base':
## 
##     date, intersect, setdiff, union
\end{verbatim}

\begin{Shaded}
\begin{Highlighting}[]
\FunctionTok{library}\NormalTok{(dplyr)}
\end{Highlighting}
\end{Shaded}

\begin{verbatim}
## 
## Attaching package: 'dplyr'
\end{verbatim}

\begin{verbatim}
## The following objects are masked from 'package:stats':
## 
##     filter, lag
\end{verbatim}

\begin{verbatim}
## The following objects are masked from 'package:base':
## 
##     intersect, setdiff, setequal, union
\end{verbatim}

\begin{Shaded}
\begin{Highlighting}[]
\FunctionTok{library}\NormalTok{(readxl)}

\CommentTok{\#Define file name and url}
\NormalTok{url }\OtherTok{\textless{}{-}} \StringTok{"https://d396qusza40orc.cloudfront.net/repdata\%2Fdata\%2Factivity.zip"}
\NormalTok{temp\_zip }\OtherTok{\textless{}{-}} \FunctionTok{tempfile}\NormalTok{(}\AttributeTok{fileext=}\StringTok{".zip"}\NormalTok{)}

\CommentTok{\#Download the file and unzip file}
\FunctionTok{download.file}\NormalTok{(url,temp\_zip,}\AttributeTok{mode=}\StringTok{"wb"}\NormalTok{)}
\NormalTok{file\_list }\OtherTok{\textless{}{-}} \FunctionTok{unzip}\NormalTok{(temp\_zip,}\AttributeTok{list=}\ConstantTok{TRUE}\NormalTok{)}
\FunctionTok{print}\NormalTok{(file\_list) }\CommentTok{\# To see file name}
\end{Highlighting}
\end{Shaded}

\begin{verbatim}
##           Name Length                Date
## 1 activity.csv 350829 2014-02-11 10:08:00
\end{verbatim}

\begin{Shaded}
\begin{Highlighting}[]
\NormalTok{activitydata }\OtherTok{\textless{}{-}} \FunctionTok{read.csv}\NormalTok{(}\FunctionTok{unz}\NormalTok{(temp\_zip,}\StringTok{"activity.csv"}\NormalTok{))}
\FunctionTok{unlink}\NormalTok{(temp\_zip)}
\FunctionTok{nrow}\NormalTok{(activitydata)}
\end{Highlighting}
\end{Shaded}

\begin{verbatim}
## [1] 17568
\end{verbatim}

\begin{Shaded}
\begin{Highlighting}[]
\FunctionTok{names}\NormalTok{(activitydata)}
\end{Highlighting}
\end{Shaded}

\begin{verbatim}
## [1] "steps"    "date"     "interval"
\end{verbatim}

After downloading the data, we will next look at the first question in
what is the mean total number of steps taken per day with the following
analyses: 1. Calculate total number of steps taken per day 2. Make a
histogram of the total number of steps taken each day 3. Calculate the
mean and median of total number of steps taken per day

Solution: To tackle the above, we will first need to convert the date
into a format that can be analysed using as.Date function. Following
which, we will need to group the number of steps per day and sum up the
data for each day. This can be done in creating another
``totalsteps\_perday'' dataset using group\_by date function and then
pipe it to summarise using summarize function to sum up the steps taken
in another variable within the dataset called ``total\_steps''. Latter
will denote the total steps in a day and each row in the dataset is one
day. We will also use the na.rm = TRUE to ignore missing values. Next,
use the function hist() on totalsteps\_perday\$total\_steps to create
the histogram.

\begin{Shaded}
\begin{Highlighting}[]
\NormalTok{knitr}\SpecialCharTok{::}\NormalTok{opts\_chunk}\SpecialCharTok{$}\FunctionTok{set}\NormalTok{(}\AttributeTok{echo =} \ConstantTok{TRUE}\NormalTok{,}\AttributeTok{warning =} \ConstantTok{FALSE}\NormalTok{, }\AttributeTok{message =} \ConstantTok{FALSE}\NormalTok{)}

\CommentTok{\#Convert Date format}
\NormalTok{activitydata}\SpecialCharTok{$}\NormalTok{date }\OtherTok{\textless{}{-}} \FunctionTok{as.Date}\NormalTok{(activitydata}\SpecialCharTok{$}\NormalTok{date,}\AttributeTok{format=}\StringTok{"\%Y{-}\%m{-}\%d"}\NormalTok{)}

\CommentTok{\#Calculate the total steps per day by using group\_by date and then use summarise to sum up the total steps for each day}
\NormalTok{totalsteps\_perday }\OtherTok{\textless{}{-}}\NormalTok{ activitydata }\SpecialCharTok{\%\textgreater{}\%} 
                        \FunctionTok{group\_by}\NormalTok{(date) }\SpecialCharTok{\%\textgreater{}\%} \FunctionTok{summarize}\NormalTok{(}\AttributeTok{total\_steps =} \FunctionTok{sum}\NormalTok{(steps,}\AttributeTok{na.rm=}\ConstantTok{TRUE}\NormalTok{))}

\CommentTok{\#Plot Histogram}
\FunctionTok{hist}\NormalTok{(totalsteps\_perday}\SpecialCharTok{$}\NormalTok{total\_steps,}\AttributeTok{main=}\StringTok{"Total Steps Taken Each Day"}\NormalTok{,}
        \AttributeTok{xlab=}\StringTok{"Steps"}\NormalTok{,}
        \AttributeTok{col=}\StringTok{"steelblue"}\NormalTok{,}
        \AttributeTok{breaks=}\DecValTok{20}\NormalTok{, }\AttributeTok{na.rm=}\ConstantTok{TRUE}\NormalTok{)}
\end{Highlighting}
\end{Shaded}

\pandocbounded{\includegraphics[keepaspectratio]{PA1_template_files/figure-html/question 1-1.png}}

\begin{Shaded}
\begin{Highlighting}[]
\CommentTok{\# Compute Mean and Median}
\NormalTok{mean\_stepsperday }\OtherTok{\textless{}{-}} \FunctionTok{mean}\NormalTok{(totalsteps\_perday}\SpecialCharTok{$}\NormalTok{total\_steps,}\AttributeTok{na.rm=}\ConstantTok{TRUE}\NormalTok{)}
\NormalTok{median\_stepsperday }\OtherTok{\textless{}{-}} \FunctionTok{median}\NormalTok{(totalsteps\_perday}\SpecialCharTok{$}\NormalTok{total\_steps,}\AttributeTok{na.rm=}\ConstantTok{TRUE}\NormalTok{)}
\FunctionTok{print}\NormalTok{(mean\_stepsperday)}
\end{Highlighting}
\end{Shaded}

\begin{verbatim}
## [1] 9354.23
\end{verbatim}

\begin{Shaded}
\begin{Highlighting}[]
\FunctionTok{print}\NormalTok{(median\_stepsperday)}
\end{Highlighting}
\end{Shaded}

\begin{verbatim}
## [1] 10395
\end{verbatim}

Next, we will study the data to derive the answers to the below: 1) What
is the average daily activity pattern? 2) Make a time series plot of the
5-minute interval (x-axis) and the average number of steps taken,
averaged across all days (y-axis) 3) Which 5-minute interval, on average
across all the days in the dataset, contains the maximum number of
steps?

Solution: To solve the above, we will also use the group\_by and
summarize functions to create a dataset ``Average\_daily\_pattern'' that
groups by interval and then take the mean of the steps within the
interval to create a new variable ``avg\_steps'' for this. Similarly,
remove NA values by setting na.rm = TRUE.Then plot interval and
avg\_steps in the ``Average\_daily\_pattern'' dataset.To calculate the
maximum steps, we use the filter function to detect the datapoint
whereby avg\_steps is the max(avg\_steps) and then come back with the
corresponding interval which is 835.

\begin{Shaded}
\begin{Highlighting}[]
\NormalTok{knitr}\SpecialCharTok{::}\NormalTok{opts\_chunk}\SpecialCharTok{$}\FunctionTok{set}\NormalTok{(}\AttributeTok{echo =} \ConstantTok{TRUE}\NormalTok{,}\AttributeTok{warning =} \ConstantTok{FALSE}\NormalTok{, }\AttributeTok{message =} \ConstantTok{FALSE}\NormalTok{)}

\CommentTok{\# Average steps per 5{-}minute interval}
\NormalTok{Average\_daily\_pattern }\OtherTok{\textless{}{-}}\NormalTok{ activitydata }\SpecialCharTok{\%\textgreater{}\%}
                                \FunctionTok{group\_by}\NormalTok{(interval) }\SpecialCharTok{\%\textgreater{}\%}
                                \FunctionTok{summarize}\NormalTok{(}\AttributeTok{avg\_steps =} \FunctionTok{mean}\NormalTok{(steps,}\AttributeTok{na.rm=}\ConstantTok{TRUE}\NormalTok{))}
\CommentTok{\#Time series plot}
\FunctionTok{plot}\NormalTok{(Average\_daily\_pattern}\SpecialCharTok{$}\NormalTok{interval,}
\NormalTok{     Average\_daily\_pattern}\SpecialCharTok{$}\NormalTok{avg\_steps,}
        \AttributeTok{type =} \StringTok{"l"}\NormalTok{, }\AttributeTok{col=}\StringTok{"red"}\NormalTok{,}\AttributeTok{iwd =} \DecValTok{2}\NormalTok{, }\AttributeTok{main =} \StringTok{"Average Daily Activity Pattern By 5{-}minute Intervals"}\NormalTok{,}
     \AttributeTok{xlab =} \StringTok{"5{-}minute Interval"}\NormalTok{,}
     \AttributeTok{ylab =} \StringTok{"Average Steps"}\NormalTok{)}
\end{Highlighting}
\end{Shaded}

\pandocbounded{\includegraphics[keepaspectratio]{PA1_template_files/figure-html/question 2-1.png}}

\begin{Shaded}
\begin{Highlighting}[]
\CommentTok{\# Find 5{-}minute interval with maximum number of steps}
\NormalTok{max\_interval }\OtherTok{\textless{}{-}}\NormalTok{ Average\_daily\_pattern }\SpecialCharTok{\%\textgreater{}\%} \FunctionTok{filter}\NormalTok{(avg\_steps}\SpecialCharTok{==}\FunctionTok{max}\NormalTok{(avg\_steps))}
\FunctionTok{print}\NormalTok{(max\_interval}\SpecialCharTok{$}\NormalTok{interval)}
\end{Highlighting}
\end{Shaded}

\begin{verbatim}
## [1] 835
\end{verbatim}

Imputing missing values Note that there are a number of days/intervals
where there are missing values (coded as NA). The presence of missing
days may introduce bias into some calculations or summaries of the data.

\begin{enumerate}
\def\labelenumi{\arabic{enumi}.}
\item
  Calculate and report the total number of missing values in the dataset
  (i.e.~the total number of rows with NAs)
\item
  Devise a strategy for filling in all of the missing values in the
  dataset. For example, you could use the mean/median for that day, or
  the mean for that 5-minute interval, etc.
\item
  Create a new dataset that is equal to the original dataset but with
  the missing data filled in.
\item
  Make a histogram of the total number of steps taken each day and
  Calculate and report the mean and median total number of steps taken
  per day. Do these values differ from the estimates from the first part
  of the assignment? What is the impact of imputing missing data on the
  estimates of the total daily number of steps?
\end{enumerate}

Solution: To resolve the above, we use sum(is.na) to count the total
NAs. Following which, we will create a new variable
``activity\_imputed'' that imputes the missing values by first joining
the raw activity data with average daily pattern and then use mutate to
create a new variable using ifelse that takes the value of
``avg\_steps'' if the value is NA or missing. Following which, we remove
the avg\_step field from the dataset to keep it clean. With this
dataset, we recompute the daily steps using group by and summarise
functions. And then plot the histogram, mean and median.

\begin{Shaded}
\begin{Highlighting}[]
\NormalTok{knitr}\SpecialCharTok{::}\NormalTok{opts\_chunk}\SpecialCharTok{$}\FunctionTok{set}\NormalTok{(}\AttributeTok{echo =} \ConstantTok{TRUE}\NormalTok{,}\AttributeTok{warning =} \ConstantTok{FALSE}\NormalTok{, }\AttributeTok{message =} \ConstantTok{FALSE}\NormalTok{)}
\CommentTok{\# Count total NAs}
\NormalTok{total\_na }\OtherTok{\textless{}{-}} \FunctionTok{sum}\NormalTok{(}\FunctionTok{is.na}\NormalTok{(activitydata}\SpecialCharTok{$}\NormalTok{steps))}
\FunctionTok{paste0}\NormalTok{(}\StringTok{"The number of missing values are "}\NormalTok{,total\_na,}\StringTok{" ."}\NormalTok{)}
\end{Highlighting}
\end{Shaded}

\begin{verbatim}
## [1] "The number of missing values are 2304 ."
\end{verbatim}

\begin{Shaded}
\begin{Highlighting}[]
\CommentTok{\#Create a new dataset and fill NAs as average of steps }
\NormalTok{activity\_imputed }\OtherTok{\textless{}{-}}\NormalTok{ activitydata }\SpecialCharTok{\%\textgreater{}\%}
                        \FunctionTok{left\_join}\NormalTok{(Average\_daily\_pattern, }\AttributeTok{by=}\StringTok{"interval"}\NormalTok{) }\SpecialCharTok{\%\textgreater{}\%}
                        \FunctionTok{mutate}\NormalTok{(}\AttributeTok{steps =} \FunctionTok{ifelse}\NormalTok{(}\FunctionTok{is.na}\NormalTok{(steps),avg\_steps,steps)) }\SpecialCharTok{\%\textgreater{}\%}
                        \FunctionTok{select}\NormalTok{ (}\SpecialCharTok{{-}}\NormalTok{avg\_steps)}

\CommentTok{\#Recompute total steps per day with imputed data}
\NormalTok{imputed\_steps\_per\_day }\OtherTok{\textless{}{-}}\NormalTok{ activity\_imputed }\SpecialCharTok{\%\textgreater{}\%} 
                        \FunctionTok{group\_by}\NormalTok{(date) }\SpecialCharTok{\%\textgreater{}\%}
                        \FunctionTok{summarize}\NormalTok{(}\AttributeTok{total\_steps =} \FunctionTok{sum}\NormalTok{(steps,}\AttributeTok{na.rm=}\ConstantTok{TRUE}\NormalTok{))}

\CommentTok{\#Histogram for imputed data}
\FunctionTok{hist}\NormalTok{(imputed\_steps\_per\_day}\SpecialCharTok{$}\NormalTok{total\_steps,}
     \AttributeTok{main=}\StringTok{"Total Steps Per Day Using Imputed Data for NAs"}\NormalTok{,}
     \AttributeTok{xlab =} \StringTok{"Steps"}\NormalTok{,}
     \AttributeTok{col =} \StringTok{"lightgreen"}\NormalTok{,}
     \AttributeTok{breaks =} \DecValTok{20}\NormalTok{)}
\end{Highlighting}
\end{Shaded}

\pandocbounded{\includegraphics[keepaspectratio]{PA1_template_files/figure-html/question 3-1.png}}

\begin{Shaded}
\begin{Highlighting}[]
\CommentTok{\# Compute Mean and Median of imputed data}
\NormalTok{mean\_imputed\_stepsperday }\OtherTok{\textless{}{-}} 
        \FunctionTok{mean}\NormalTok{(imputed\_steps\_per\_day}\SpecialCharTok{$}\NormalTok{total\_steps,}\AttributeTok{na.rm=}\ConstantTok{TRUE}\NormalTok{)}
\NormalTok{median\_imputed\_stepsperday }\OtherTok{\textless{}{-}} 
        \FunctionTok{median}\NormalTok{(imputed\_steps\_per\_day}\SpecialCharTok{$}\NormalTok{total\_steps,}\AttributeTok{na.rm=}\ConstantTok{TRUE}\NormalTok{)}
\FunctionTok{print}\NormalTok{(mean\_imputed\_stepsperday)}
\end{Highlighting}
\end{Shaded}

\begin{verbatim}
## [1] 10766.19
\end{verbatim}

\begin{Shaded}
\begin{Highlighting}[]
\FunctionTok{print}\NormalTok{(median\_imputed\_stepsperday)}
\end{Highlighting}
\end{Shaded}

\begin{verbatim}
## [1] 10766.19
\end{verbatim}

Are there differences in activity patterns between weekdays and
weekends? For this part the weekdays() function may be of some help
here. Use the dataset with the filled-in missing values for this part.

\begin{enumerate}
\def\labelenumi{\arabic{enumi}.}
\tightlist
\item
  Create a new factor variable in the dataset with two levels b
\end{enumerate}

\end{document}
